% ===== Block B01: 原文 p.1–17 | §1.1 紧黎曼曲面 | 译手 tr-ch1a =====
\subsection{紧黎曼曲面}
% [原文 p.1]
\noindent\textbf{1.1.1 紧黎曼曲面.}\quad 一个连通的紧一维复流形称为紧黎曼曲面（compact Riemann surface），或闭黎曼曲面。在本讲义中我们主要考虑紧黎曼曲面；我们常常省略「紧」字。

紧黎曼曲面 $R$ 是一个可定向的闭曲面，因此微分同胚于有 $g$ 个洞的甜甜圈曲面。数 $g$ 称为黎曼曲面 $R$ 的亏格（genus），常记为 $g(R)$。

对亏格为 $g$ 的紧黎曼曲面 $R$，其一阶同调群 $H_1(R,\Z)$ 是秩为 $2g$ 的自由 $\Z$-模。此外，相交配对 $\alpha\cdot\beta$（$\alpha,\beta\in H_1(R,\Z)$）是斜对称的，并且给出双线性映射
\[
  H_1(R,\Z)\times H_1(R,\Z)\longrightarrow \Z .
\]

我们总能找到 $H_1(R,\Z)$ 的一组辛基（symplectic basis）$\{\alpha_1,\dots,\alpha_g,\beta_1,\dots,\beta_g\}$，即
\[
  \alpha_i\cdot\alpha_j=0,\qquad \beta_i\cdot\beta_j=0,\qquad \alpha_i\cdot\beta_j=\delta_{ij}.
\]

此外，$R$ 上有 $g$ 个线性无关的全纯 1-形式。设 $\Theta_R$ 是 $R$ 上全纯向量场芽的层，$\omega_R$ 是它的对偶层。层 $\omega_R$ 称为典范层（canonical sheaf）。它是与一个复线丛 $K_X$ 相伴的层，而这个线丛称为 $R$ 的典范线丛（canonical line bundle）。众所周知，若 $R$ 上线丛 $L$ 的次数至少为 $2g+1$，则 $H^0(R,\cO_R(L))$ 给出 $R$ 到复射影空间 $\PP^{2g-3}$ 的全纯嵌入。因此所有紧黎曼曲面都具有射影簇的结构。以下我们有时把紧黎曼曲面看作射影簇。

% [Tutor Note: 笔误：$K_X$ 的下标应为 $R$——同句已说 "of $R$"，且全节只有曲面 $R$
%   一个记号，$X$ 无来源。照原文排，不改。原文 p.1
%  （另）疑误：$\PP^{2g-3}$ 的维数不可能对。$H^0(R,\cO_R(L))$ 在
%   $\deg L\ge 2g-1$ 时满足 $\dim = \deg L-g+1$（Riemann–Roch），故嵌入的靶空间
%   应为 $\PP^{\deg L-g}$；当 $\deg L=2g+1$ 时是 $\PP^{g+1}$。核对方法：$g=1$ 时
%   $2g-3=-1$ 显然荒谬；$g=2$、$L=K$ 时 $\deg=2$ 也不能嵌入 $\PP^{1}$。原文 p.1]

\figomit{1.1}{亏格为 $g$ 的黎曼曲面。}

% [原文 p.2]
\noindent\textbf{1.1.2 黎曼曲面的复解析族.}\quad 以下，对复流形 $M$ 及其上一点 $Q$，我们用 $T_Q(M)$ 表示 $Q$ 处的全纯切向量空间。此外，除非特别声明，我们总假定复流形是连通的。

设 $\pi:\mathcal{C}\to\cW$ 是从 $(m+1)$ 维复流形 $\mathcal{C}$ 到 $m$ 维复流形 $\cW$ 的全纯映射，若它满足下列条件，则称为 $\cW$ 上紧黎曼曲面的光滑族（smooth family of compact Riemann surfaces），或紧黎曼曲面的复解析族（complex analytic family of compact Riemann surfaces）：
\begin{enumerate}[label=(\arabic*)]
\item 映射 $\pi$ 是逆紧的（proper）。即：对复流形 $\cW$ 的任意紧集 $K$，原像 $\pi^{-1}(K)$ 是紧的。
\item 映射 $\pi$ 是光滑且全纯的。即：对任意点 $P\in\mathcal{C}$，全纯切空间之间的线性映射 $(d\pi)_P:T_P(\mathcal{C})\to T_{\pi(P)}(\cW)$ 是满射。
\item 对任意点 $w\in\cW$，纤维 $\pi^{-1}(w)$ 是连通的。
\end{enumerate}

由条件 (1) 与 (2)，每一点 $w\in\cW$ 上方的纤维 $C_w=\pi^{-1}(w)$ 是紧黎曼曲面。又，对一点 $w_0\in\cW$，我们把 $C_w$（$w\in\cW$）称为紧黎曼曲面 $C_{w_0}$ 的一个形变（deformation）。

当给定紧黎曼曲面的复解析族 $\pi:\mathcal{C}\to\cW$ 与全纯映射 $h:\mathcal{S}\to\cW$ 时，纤维积 $\mathcal{C}\times_\cW\mathcal{S}\to\mathcal{S}$ 是 $\mathcal{S}$ 上的复解析族。我们把族 $\mathcal{C}\times_\cW\mathcal{S}\to\mathcal{S}$ 称为由全纯映射 $h$ 拉回的拉回族（pullback family）。

% [Tutor Note: 术语：这里 "complex analytic family" 与 "smooth family" 是同义语，
%  glossary 定为「复解析族」；"deformation" 定为「形变」，注意它在第 3 节
%  （§1.3 Deformation of pointed curves）才系统展开，此处只是命名。原文 p.2]

\begin{example}[1.1]
在上半平面 $H=\{\tau\in\C\mid \operatorname{Im}\tau>0\}$ 与复平面 $\C$ 的乘积 $H\times\C$ 上，群 $\Z^2$ 的作用定义为
\[
\begin{aligned}
(m,n):\quad & H\times\C \longrightarrow H\times\C,\qquad (m,n)\in\Z^2,\\
& (\tau,\zeta)\longmapsto (\tau,\zeta+m+n\tau).
\end{aligned}
\]
由于该作用无不动点且真不连续，商空间 $\cE=H\times\C/\Z^2$ 具有二维复流形的结构。点 $(\tau,\zeta)\in H\times\C$ 决定 $\cE$ 的点 $(\tau,[\zeta])$。自然投影 $p:H\times\C\to H$ 诱导全纯映射
\[
\begin{aligned}
\pi:\quad & \cE \longrightarrow H,\\
& (\tau,[\zeta])\longmapsto \tau .
\end{aligned}
\]
全纯映射 $\pi:\cE\to H$ 满足上面的条件 (1)、(2)、(3)，因此是 $H$ 上紧黎曼曲面的复解析族。对点 $\tau\in H$，纤维 $E_\tau=\pi^{-1}(\tau)$ 是以 $(1,\tau)$ 为基本周期的椭圆曲线（一维复环面）。
\end{example}

\begin{example}[1.2]
群 $SL(2,\R)$ 通过
\[
\tau\longmapsto \frac{a\tau+b}{c\tau+d},\qquad \begin{pmatrix} a & b\\ c & d\end{pmatrix}\in SL(2,\R)
\]
作用在上半平面 $H$ 上。

对正整数 $n$，$SL(2,\Z)$ 的正规子群 $\Gamma(n)$ 定义为
\[
\Gamma(n)=\left\{\begin{pmatrix} a & b\\ c & d\end{pmatrix}\in SL(2,\Z)\ \middle|\
a\equiv d\equiv 1\ (\mathrm{mod}\ n),\ b\equiv c\equiv 0\ (\mathrm{mod}\ n)\right\}.
\]
于是众所周知，当 $n\ge 3$ 时群 $\Gamma(n)$ 在 $H$ 上的作用没有不动点。此外，群 $\Gamma(n)$ 通过
\[
(\tau,[\zeta])\longmapsto \left(\frac{a\tau+b}{c\tau+d},\ \left[\frac{\zeta}{c\tau+d}\right]\right)
\]
作用于例 1.1 的 $\cE$ 上。

% [原文 p.3]
该作用无不动点。全纯映射 $\pi:\cE\to H$ 与 $\Gamma(n)$ 在 $\cE$ 与 $H$ 上的作用相容。因此我们得到商空间之间的全纯映射 $\pi(n):\cE(n):=\cE/\Gamma(n)\to C(n):=H/\Gamma(n)$。

对 $n\ge3$，$\pi(n):\cE(n)\to C(n)$ 也是紧黎曼曲面的复解析族。
\end{example}

% [Tutor Note: $\Gamma(n)$ 的作用与例 1.1 的 $\Z^2$ 作用可交换，才可能在
%   $\cE=H\times\C/\Z^2$ 上诱导作用；核对：$SL(2,\Z)$ 把格 $\Z+\Z\tau$ 送到
%   $\Z+\Z\frac{a\tau+b}{c\tau+d}$，而 $\Gamma(n)$ 在 $\tau$ 上的作用与 $\zeta$ 的
%   变换 $\zeta/(c\tau+d)$ 一起保持 $\Z^2$-轨道。原文 p.2
%  （另）记号提醒：$C(n)=H/\Gamma(n)$ 是模曲线，$n\ge 3$ 时 $\Gamma(n)$ 无椭圆元、
%   无抛物元，故 $C(n)$ 光滑而 $\cE(n)\to C(n)$ 是椭圆曲线的通用族；$n=1,2$ 时
%   $C(n)$ 有椭圆型与尖点，族不再是"紧黎曼曲面的族"。原文 p.3]

\noindent\textbf{1.1.3 小平–Spencer 映射.}\quad 让我们用局部坐标来研究紧黎曼曲面的复解析族 $\pi:\mathcal{C}\to\cW$。取 $\cW$ 中点 $0$ 的一个坐标邻域 $U$，以及 $U$ 上以 $0$ 为中心的局部坐标 $(w^1,\dots,w^m)$。设 $\pi^{-1}(U)$ 被开集 $\{U_\lambda\}_{\lambda\in\Lambda}$ 覆盖，即 $\bigcup_{\lambda\in\Lambda}U_\lambda=\pi^{-1}(U)$。由于 $\pi$ 是光滑全纯映射，我们可以取 $U_\lambda$ 的局部坐标为 $(w^1,\dots,w^m,z_\lambda)$。（也就是说，对每个 $\lambda$，$w^1,\dots,w^m$ 都被取作 $U_\lambda$ 的局部坐标的一部分。）若 $U_\lambda\cap U_\mu\ne\varnothing$，则在两套局部坐标之间有下述关系：
\begin{equation}
  z_\lambda=f_{\lambda\mu}(z_\mu,w^1,\dots,w^m). \tag{1.1}
\end{equation}

这里 $f_{\lambda\mu}$ 在 $U_\lambda\cap U_\mu$ 上全纯。我们可以把复流形 $\pi^{-1}(U)$ 看作通过关系式 (1.1) 粘合 $\{U_\lambda\}_{\lambda\in\Lambda}$ 而得到的。于是 $(w^1,\dots,w^m)$ 可以看作改变粘合方式的参数。也就是说，对 $(w^1,\dots,w^m)=(0,\dots,0)$，流形 $\pi^{-1}((0,\dots,0))$ 是紧黎曼曲面 $C_0$，它由 $z_\lambda=f_{\lambda\mu}(z_\mu,0,\dots,0)$ 粘合得到。而对 $w=(w^1,\dots,w^m)$，我们通过对 $C_0$ 的粘合方式作关于 $w$ 的微小改变而得到 $C_w=\pi^{-1}(w)$。因此，(1.1) 关于变量 $w^1,\dots,w^m$ 的 Taylor 展开的一次项给出复流形 $C_0$ 的形变的一阶近似（通常称为无穷小形变（infinitesimal deformation)）。这意味着，若 $U_\lambda\cap U_\mu\ne\varnothing$，则
\[
\left\{\frac{\partial f_{\lambda\mu}}{\partial w^k}(z_\mu,0,\dots,0)\right\},\qquad 1\le k\le m,
\]
给出了关于紧黎曼曲面 $C_0$ 的形变的信息。

对 $U_\lambda\cap U_\mu\ne\varnothing$，让我们考虑全纯切向量场
\begin{equation}
  \theta^{(k)}_{\lambda\mu}=\frac{\partial f_{\lambda\mu}}{\partial w^k}(z_\mu,0,\dots,0)\frac{\partial}{\partial z_\lambda} \tag{1.2}
\end{equation}
（定义在 $C_0\cap U_\lambda\cap U_\mu$ 上）。注意在上述向量场 $\theta^{(k)}_{\lambda\mu}$ 的定义中，我们用的是 $U_\lambda$ 上的向量场 $\frac{\partial}{\partial z_\lambda}$，而不是 $\frac{\partial}{\partial z_\mu}$。若有 $U_\lambda\cap U_\mu\cap U_\nu\ne\varnothing$，则由 (1.1) 我们有
\[
\begin{aligned}
z_\lambda &= f_{\lambda\mu}\left(f^{1}_{\mu\nu}(z_\nu,w),w^1,\dots,w^m\right)\\
&= f_{\lambda\nu}(z_\nu,w^1,\dots,w^m).
\end{aligned}
\]

% [Tutor Note: $f^{1}_{\mu\nu}$ 表示 $f_{\mu\nu}$ 关于第一个变量 $z_\nu$ 的反函数
%  （即 $z_\nu\mapsto z_\mu$ 的逆映射），不是 $1/f_{\mu\nu}$；这里的上标是"反函数"
%  的旧记号，见下一行把 $z_\mu=f_{\mu\nu}(z_\nu,w)$ 代入左边即得。原文 p.3]
% [Tutor Note: (1.2) 里必须用 $\partial/\partial z_\lambda$ 而不能用 $\partial/\partial z_\mu$：
%  因为 $\theta^{(k)}_{\lambda\mu}$ 要在 $U_\lambda$ 的坐标卡里是一个向量场（切于纤维），
%  换到 $U_\mu$ 的坐标会差一个 Jacobi 因子 $\partial z_\lambda/\partial z_\mu$，
%  那样就不再是同一个上链的值。这个"用哪一边"的选择是后面 (1.3) 良定义性的来源。原文 p.3]

% [原文 p.4]
于是，在 $U_\lambda\cap U_\mu\cap U_\nu$ 上有
\[
\begin{aligned}
\theta^{(k)}_{\lambda\nu}
&= \frac{\partial f_{\lambda\nu}}{\partial w^k}(z_\nu,0)\frac{\partial}{\partial z_\lambda}\\
&= \frac{\partial f_{\lambda\mu}}{\partial z_\mu}(f_{\mu\nu}(z_\nu,0),0)\cdot\frac{\partial f_{\mu\nu}}{\partial w^k}(z_\nu,0)\frac{\partial}{\partial z_\lambda}\\
&\quad +\frac{\partial f_{\lambda\mu}}{\partial w^k}(f_{\mu\nu}(z_\nu,0),0)\frac{\partial}{\partial z_\lambda}\\
&= \frac{\partial f_{\mu\nu}}{\partial w^k}(z_\nu,0)\left(\frac{\partial f_{\lambda\mu}}{\partial z_\mu}\frac{\partial}{\partial z_\lambda}\right)+\theta^{(k)}_{\lambda\mu}\\
&= \frac{\partial f_{\mu\nu}}{\partial w^k}(z_\nu,0)\frac{\partial}{\partial z_\mu}+\theta^{(k)}_{\lambda\mu}\\
&= \theta^{(k)}_{\mu\nu}+\theta^{(k)}_{\lambda\mu}.
\end{aligned}
\]

也就是说，$\{\theta^{(k)}_{\lambda\mu}\}$ 是系数取值于全纯切向量场的 Čech 一上链。由上述一上链定义的上同调类 $H^1(C_0,\Theta)$ 中的元素仍用同一记号 $\{\theta^{(k)}_{\lambda\mu}\}$ 表示。这里，我们用 $\Theta$ 表示 $C_0$ 上全纯向量场芽的层。

% [Tutor Note: 上链关系的核对：把 $z_\lambda=f_{\lambda\mu}(f_{\mu\nu}(z_\nu,w),w)$ 对 $w^k$
%  求导，就是链式法则给出的两行之和（第一行对第二个变量求导、第二行对第一个变量
%  求导）；注意第三行用的恒等式 $\partial f_{\lambda\mu}/\partial z_\mu\cdot\partial/\partial z_\lambda
%  =\partial/\partial z_\mu$ 正是坐标变换的链式法则。若把 $\lambda,\mu$ 次序写反就会得到
%  $\theta_{\lambda\nu}=\theta_{\lambda\mu}+\theta_{\mu\nu}$ 的错误符号顺序。原文 p.4]

若固定局部坐标 $(w^1,\dots,w^m)$，则上同调类 $\{\theta^{(k)}_{\lambda\mu}\}$ 被唯一确定，并且不依赖于 $\pi^{-1}(U)$ 的开覆盖 $\{U_\lambda\}_{\lambda\in\Lambda}$ 与坐标邻域、局部坐标 $(z_\lambda,w^1,\dots,w^m)$ 的选取。这一点可由直接计算证明。因此，我们可以定义从 $0\in\cW$ 处的全纯切向量空间 $T_0\cW$ 到 $H^1(C_0,\Theta)$ 的线性映射：
\begin{equation}
\begin{aligned}
\rho_0:\quad & T_0\cW \longrightarrow H^1(C_0,\Theta),\\
& \sum a_k\frac{\partial}{\partial w^k} \longmapsto \sum a_k\left\{\theta^{(k)}_{\lambda\nu}\right\}.
\end{aligned} \tag{1.3}
\end{equation}
这个线性映射称为小平–Spencer 映射（Kodaira-Spencer mapping）。

顺便说明一下，在 (1.2) 中 $\{\theta^{(k)}_{\mu\nu}\}$ 的定义里，如果我们不取 $w=0$，而取
\[
\theta^{(k)}_{\lambda\mu}(w)=\frac{\partial f_{\lambda\mu}}{\partial w^k}(z_\mu,w)\frac{\partial}{\partial z_\lambda},
\]
那么由与上面相同的论证可以证明 $\{\theta^{(k)}_{\lambda\mu}(w)\}\in H^1(C_w,\Theta_{C_w})$。此外，$\theta^{(k)}_{\lambda\mu}(w)$ 关于 $w$ 是全纯的。因此，与 (1.3) 类似，我们得到 $\cO_\cW$-同态
\begin{equation}
\begin{aligned}
\rho:\quad & \Theta_\cW \longrightarrow R^1\pi_*\Theta_{\mathcal{C}/\cW},\\
& \sum a_k(w)\frac{\partial}{\partial w^k}\longmapsto \sum a_k(w)\left\{\theta^{(k)}_{\lambda\mu}(w)\right\};
\end{aligned} \tag{1.4}
\end{equation}
这里 $\Theta_{\mathcal{C}/\cW}$ 是 $\pi:\mathcal{C}\to\cW$ 的相对全纯切向量场芽的层。$\cO_\cW$-同态 $\rho$ 也称为小平–Spencer 映射。

由于复解析族 $\pi:\mathcal{C}\to\cW$ 是光滑全纯映射，存在 $\cO_\cW$-同态的正合列
\begin{equation}
0\longrightarrow \Theta_{\mathcal{C}/\cW}\longrightarrow \Theta_{\mathcal{C}}\longrightarrow \pi^*\Theta_\cW\longrightarrow 0. \tag{1.5}
\end{equation}
由该正合列得到长正合列
\[
\longrightarrow \pi_*\Theta_{\mathcal{C}}\longrightarrow \pi_*\pi^*\Theta_\cW=\Theta_\cW\longrightarrow R^1\pi_*\Theta_{\mathcal{C}/\cW}\longrightarrow R^1\pi_*\Theta_{\mathcal{C}}\longrightarrow \cdots
\]

% [Tutor Note: 笔误：本段说"在 (1.2) 中 $\{\theta^{(k)}_{\mu\nu}\}$ 的定义里"，但 (1.2) 定义的是
%   $\theta^{(k)}_{\lambda\mu}$，全文再无 $\theta^{(k)}_{\mu\nu}$ 的定义，下标应为 $\lambda\mu$。
%   照原文排，不改。原文 p.4
%  （另）$R^1\pi_*$ 是一阶直像层（相对上同调的层化），故 (1.4) 的靶是 $\cW$ 上的层，
%   而 (1.3) 的靶只是 $0$ 点处的一个向量空间；二者关系是 $\rho$ 在 $0$ 点的纤维给出 $\rho_0$。
%   注意 (1.4) 中 $\Theta_\cW$ 是 $\cW$ 的切层，$\Theta_{\mathcal{C}/\cW}$ 是相对切层。原文 p.4
%  （另）(1.5) 的正合性只用 $\pi$ 光滑：$d\pi:\Theta_{\mathcal{C}}\to\pi^*\Theta_\cW$ 满射，
%   其核就是相对切层 $\Theta_{\mathcal{C}/\cW}$。（1.5）没有用到纤维连通性；连通性是在
%   保证纤维为紧黎曼曲面时用的。原文 p.4]
% [原文 p.5]
由 $\rho$ 的定义容易看出，小平–Spencer 映射 (1.4) 正是上述长正合列中的 $\cO_\cW$-同态 $\Theta_\cW\to R^1\pi_*\Theta_{\mathcal{C}/\cW}$。

\begin{remark}[1.3]
类似地，我们可以定义紧复流形的复解析族，及其小平–Spencer 映射。
\end{remark}

下面我们计算例 1.2 中构造的椭圆曲线族的小平–Spencer 映射。

% [Tutor Note: 疑误：这里与 p.5 下一段两次引用的 "Example 1.2" 都应为例 1.1——椭圆
%   曲线族 $\cE\to H$ 是例 1.1 构造的，例 1.2 只是给出 $\Gamma(n)$ 在它上面的作用。
%   照原文排，不改。原文 p.5
%  （另）这条上链的形状来自"把开集 $\cU_0$ 上的局部截面 $z^{-1}\frac{d}{dz}$ 搬到
%   其他开集上"：在 $\cU_\mu$（$\mu\ge1$）上它全纯，所以取 $0$；交错项因此只在
%   $(\lambda,\mu)=(0,\mu_{\ge1})$ 与 $(\lambda_{\ge1},0)$ 出现，符号相反以保证
%   $\delta\tau=0$（Čech 上链条件）。若两处都用 $+$，$\delta\tau\ne0$。原文 p.5]

\begin{example}[1.4]
我们先考虑椭圆曲线 $E=\C/(1,\tau_0)$，$\tau_0\in H$。设 $z$ 是仿射空间 $\C$ 的整体坐标，$o$ 是椭圆曲线 $E$ 的原点。由正合列
\[
0\longrightarrow\Theta_E\longrightarrow\Theta_E(*o)\xrightarrow{\;t\;} z^{-1}\C[z^{-1}]\frac{d}{dz}\longrightarrow 0
\]
我们得到正合列
\[
0\longrightarrow H^0(E,\Theta_E(*o))\xrightarrow{\;t\;} z^{-1}\C[z^{-1}]\frac{d}{dz}\xrightarrow{\;\theta\;} H^1(E,\Theta_E)\longrightarrow 0,
\]
其中我们置
\[
\begin{aligned}
\Theta_E(*o)&=\varprojlim_{m\to\infty}\Theta(mo),\\
H^0(E,\Theta_E(*o))&=\varprojlim_{m\to\infty}H^0(E,\Theta(mo)).
\end{aligned}
\]
注意对 Weierstrass $\wp$ 函数，我们有
\[
\frac{(-1)^n}{n!}\wp(z)\frac{d}{dz}\in H^0(E,\cO_E((n+2)o))\setminus H^0(E,\cO_E((n+1)o))
\]
对一切非负整数 $n$ 成立。因此 $z^{-1}\frac{d}{dz}$ 在 $H^1(E,\Theta_E)$ 中有非零的像 $\theta\left(z^{-1}\frac{d}{dz}\right)$。下面我们描述这个像。设 $\{\cU_\lambda\}_{\lambda=0}^N$ 是 $E$ 的一个足够小的开覆盖，使得 $\cU_0$ 是 $o$ 的坐标邻域、局部坐标为 $z$，且满足
\[
o\notin\cU_0\cap\cU_\mu,\qquad \mu\ge 1.
\]
那么 $z^{-1}\frac{d}{dz}$ 的像由上链
\[
\tau_{\lambda\mu}=
\begin{cases}
z^{-1}\dfrac{d}{dz}, & \lambda=0,\ \mu\ge1,\ \cU_{\lambda\mu}\ne\emptyset,\\[2mm]
-z^{-1}\dfrac{d}{dz}, & \lambda\ge1,\ \mu=0,\ \cU_{\lambda\mu}\ne\emptyset,\\[2mm]
0, & \text{其他}
\end{cases}
\]
给出。

设 $p:\cE\to H$ 是例 1.2 中构造的通用族。由于我们已置 $E=\C/(1,\tau_0)$，利用上述覆盖，我们得到 $p^{-1}(V)$ 的一个开覆盖 $\{\tilde{\cU}_\lambda\}_{\lambda=0}^N$，其中 $V$ 是 $\tau_0$ 的开邻域，且 $p^{-1}(\tau_0)\cap\tilde{\cU}_\lambda=\cU_\lambda$。设 $(\tau_\lambda,z_\lambda)$ 是 $\tilde{\cU}_\lambda$ 的局部坐标，使得
\[
\begin{aligned}
  \tau_\lambda &= \tau-\tau_0,\\
  z_0 &= z,\\
  z_\lambda &= z_\mu+m_{\lambda\mu}+n_{\lambda\mu}\tau
\end{aligned}
\]
% [原文 p.6]
其中 $m_{\lambda\mu},n_{\lambda\mu}\in\Q$。那么，小平–Spencer 类 $\rho_{\tau_0}\left(\frac{\partial}{\partial\tau}\right)\in H^1(E,\Theta_E)$ 由上链
\[
  \theta_{\lambda\mu}=\left\{n_{\lambda\mu}\frac{d}{dz_\lambda}\right\}
\]
给出。

让我们考虑上链 $\{\tau_{\lambda\mu}\}$ 与 $\{\theta_{\lambda\mu}\}$ 的 Dolbeault 表示。首先考虑上链 $\{\theta_{\lambda\mu}\}$。置
\[
  \psi_\lambda=\frac{(\bar z_\lambda-z_\lambda)}{2\sqrt{-1}\operatorname{Im}\tau}\frac{d}{dz_\lambda}.
\]
那么我们有
\[
  \psi_\mu-\psi_\lambda=n_{\lambda\mu}\frac{d}{dz_\lambda}=\theta_{\lambda\mu}.
\]
因此
\[
  \bar\partial\psi_\mu=\bar\partial\psi_\lambda=(2\sqrt{-1}\operatorname{Im}\tau)^{-1}\frac{d}{dz}d\bar z
\]
就是 $\{\theta_{\lambda\mu}\}$ 的相应 Dolbeault 表示。

% [Tutor Note: $\psi_\lambda=\frac{(\bar z_\lambda-z_\lambda)}{2\sqrt{-1}\operatorname{Im}\tau}\frac{d}{dz_\lambda}$
%   本身是 $(1,0)$-形式（$\frac{d}{dz}$ 型），但不是张量（换坐标时多出的项由 $\bar z_\lambda$ 的变换抵消）；
%   关键在于 $\psi_\mu-\psi_\lambda$ 把差的 $(0,0)$ 部分抵消成
%   $n_{\lambda\mu}\frac{d}{dz_\lambda}$，而 $\bar\partial\psi_\lambda$ 与 $\lambda$ 无关，
%   故 $\{\theta_{\lambda\mu}\}$ 的 Dolbeault 代表是整体的 $(1,1)$-形式
%   $(2\sqrt{-1}\operatorname{Im}\tau)^{-1}\frac{d}{dz}d\bar z$。核对 $\bar z_\lambda=\bar z_\mu$
%   用于消去 $\bar z$ 项。原文 p.6
%  （另）$\tau_{\lambda\mu}=\tilde\tau_\mu-\tilde\tau_\lambda$ 的符号与 p.5 的上链定义
%   相反是**相容的**：Čech 上链是 $\tau_{\lambda\mu}=\tilde\tau_\mu-\tilde\tau_\lambda$，
%   而 p.5 表中 $(\lambda=0,\mu\ge1)$ 取 $+z^{-1}\frac{d}{dz}$，因为 $\tilde\tau_0$ 带
%   $f(z)/z$ 而 $\tilde\tau_\mu=0$。核对时不要照抄"$\tau=\tilde\tau_\lambda-\tilde\tau_\mu$"。
%   原文 p.6
%  （另）(1.6) 的取法要求 $f$ 在 $|z|<\epsilon/3$ 内为 $0$、在 $|z|\ge\epsilon/2$ 外为 $1$；
%   这保证 $\tilde\tau_0$ 在 $z\ne0$ 附近全纯、在 $|z|<\epsilon/3$ 上为 $0$，从而
%   $\bar\partial\tilde\tau_0$ 的支集落在环带 $\epsilon/3\le|z|\le\epsilon/2$ 内。原文 p.6]

其次考虑上同调类 $\{\tau_{\lambda\mu}\}$。我们可以假设
\[
\{z\in\C\mid |z|\le\epsilon\}\subset\cU_0
\]
以及
\[
\{z\in\C\mid |z|\le\epsilon/2\}\subset\cU_0\setminus\bigcup_{\lambda=1}^N\cU_\lambda .
\]
设 $f(z)$ 是 $\cU_0$ 上的实值 $C^\infty$ 函数，使得 $0\le f(z)\le 1$ 且
\begin{equation}
f(z)=
\begin{cases}
0, & |z|<\epsilon/3,\\
1, & |z|\ge\epsilon/2.
\end{cases} \tag{1.6}
\end{equation}
置
\[
\tilde\tau_\mu=
\begin{cases}
\dfrac{f(z)}{z}\dfrac{d}{dz}, & \mu=0,\\[2mm]
0, & \mu\ge 1.
\end{cases}
\]
那么我们有
\[
\tau_{\lambda\mu}=\tilde\tau_\mu-\tilde\tau_\lambda
\]
以及
\[
\omega=\bar\partial\tilde\tau_\lambda=\frac{1}{z}\left(\frac{\partial f(z)}{\partial\bar z}\right)\frac{d}{dz}d\bar z
\]
表示 $\{\tau_{\lambda\mu}\}$ 的 Dolbeault 上同调类。由于 Dolbeault 上同调类有形状为
\[
a\frac{d}{dz}d\bar z,\qquad a\in\C
\]
的代表，我们将对上面定义的 $\omega$ 求出常数 $a$。注意对 $E$ 上任一 $C^\infty$ 函数 $h$，我们有
\[
\int_E(\omega+\bar\partial h)\wedge dz=\int_E\omega\wedge dz .
\]

% [原文 p.7]
我们有
\[
\begin{aligned}
  \int_{E-\{|z|<\epsilon\}}\frac{\partial f(z)}{\partial\bar z}\frac{1}{z}d\bar z\wedge dz
  &= \int_{E-\{|z|<\epsilon\}}\bar\partial\left(\frac{f(z)}{z}dz\right)\\
  &= \int_{E-\{|z|<\epsilon\}}d\left(\frac{f(z)}{z}dz\right)\\
  &= -\int_{|z|=\epsilon}\frac{1}{z}dz=-2\pi\sqrt{-1}
\end{aligned}
\]
（由 (1.6)）。另一方面，我们有
\[
  \int_E a\,d\bar z\wedge dz=2\sqrt{-1}a\operatorname{Im}\tau .
\]
因此我们断言 $\{\tau_{\lambda\mu}\}$ 与 $-\frac{\pi}{\operatorname{Im}\tau}\frac{d}{dz}d\bar z$ 表示同一个上同调类。于是我们得到等式
\[
  \theta\left(z^{-1}\frac{d}{dz}\right)=-2\pi\sqrt{-1}\rho\left(\frac{\partial}{\partial\tau}\right).
\]

% [Tutor Note: 第三行用 Stokes：$E-\{|z|<\epsilon\}$ 的边界是圆周 $|z|=\epsilon$ 取**顺时针**
%   定向（区域在圆外），故 $\int_{\partial}= -\int_{|z|=\epsilon}$；被积函数
%   $\frac{f(z)}{z}dz$ 在环带上全纯（$f\equiv1$ 于 $|z|\ge\epsilon/2$），于是只剩
%   $-2\pi\sqrt{-1}$ 的留数项。若把定向记错会多一个负号，进而使 $a$ 的符号错。原文 p.7
%  （另）由两式可得 $\frac{\partial f}{\partial\bar z}$ 的积分给出 $-2\pi\sqrt{-1}$
%   而 $a\frac{d}{dz}d\bar z$ 给出 $2\sqrt{-1}a\operatorname{Im}\tau$，故
%   $a=-\frac{\pi}{\operatorname{Im}\tau}$；这一步同时验证了 $\omega$ 与
%   $-\frac{\pi}{\operatorname{Im}\tau}\frac{d}{dz}d\bar z$ 同类（差一个 $\bar\partial$ 的项）。原文 p.7]

下面我们考虑另一个坐标。置
\[
  w=\exp(2\pi\sqrt{-1}z),\qquad q_0=\exp(2\pi\sqrt{-1}\tau_0).
\]
那么椭圆曲线 $E$ 可以看作商流形 $\C^*/\langle q_0\rangle$，其中 $\langle q_0\rangle$ 表示由解析自同构
\[
  w\longmapsto q_0w
\]
生成的、作用在 $\C^*$ 上的无限循环群。在这个坐标下，椭圆曲线的原点对应于 $w=1$。置
\[
  u=w-1 .
\]
那么我们有正合列
\[
  0\longrightarrow \Theta_E\longrightarrow \Theta_E(*o)\xrightarrow{\ \eta\ } u^{-1}\C[u^{-1}]\frac{d}{du}\longrightarrow 0
\]
并由这个正合列得到正合列
\[
  0\longrightarrow H^0(E,\Theta_E(*o))\xrightarrow{\ t\ } u^{-1}\C[u^{-1}]\frac{d}{du}\xrightarrow{\ \theta\ } H^1(E,\Theta_E)\longrightarrow 0 .
\]
像 $\theta\left(u^{-1}\frac{d}{du}\right)$ 计算如下。该像由上链
\[
  \omega_{\lambda\mu}=\begin{cases}
    u^{-1}\dfrac{d}{du} & \text{若 } \lambda=0,\ \mu\ge 1,\ \cU_{0\mu}\ne 0\\[2mm]
    -u^{-1}\dfrac{d}{du} & \text{若 } \lambda\ge 1,\ \mu=0,\ \cU_{\lambda 0}\ne 0\\[2mm]
    0 & \text{其他}
  \end{cases}
\]
定义。由于我们有
\[
\begin{aligned}
  u^{-1}\frac{d}{du}
  &= \frac{1}{(\exp(2\pi\sqrt{-1}z)-1)2\pi\sqrt{-1}\exp(2\pi\sqrt{-1}z)}\frac{d}{dz}\\
  &= \left(\frac{-1}{4\pi^2z}+\text{全纯}\right)\frac{d}{dz},
\end{aligned}
\]
我们断言
\[
  \{\omega_{\lambda\mu}\}=-\frac{1}{4\pi^2}\{\tau_{\lambda\tau}\},
\]

% [Tutor Note: 笔误：上面 $\cU_{0\mu}\ne 0$、$\cU_{\lambda0}\ne 0$ 中的 $0$ 应为空集
%   $\varnothing$（与 p.5 同款上链的写法一致：$\cU_{\lambda\mu}\ne\varnothing$）。照原文排。原文 p.7
%  （另）笔误：末式下标 $\{\tau_{\lambda\tau}\}$ 应为 $\{\tau_{\lambda\mu}\}$，下标 $\tau$
%   无意义（$\tau$ 是本族参数而非覆盖指标）。照原文排，不改。原文 p.7
%  （另）第二个等号的展开：$u=e^{2\pi\sqrt{-1}z}-1\approx 2\pi\sqrt{-1}z$，
%   故 $u^{-1}\frac{d}{du}\approx\frac{1}{2\pi\sqrt{-1}z}\cdot\frac{1}{2\pi\sqrt{-1}}\frac{d}{dz}
%   =\frac{-1}{4\pi^2 z}\frac{d}{dz}$；正是这条"主部系数"把两族上链按 $-1/4\pi^2$ 联系起来。
%   原文 p.7]
% [原文 p.8]
即
\[
  \theta\left(u^{-1}\frac{d}{dz}\right)=-\frac{1}{4\pi^2}\theta\left(z^{-1}\frac{d}{dz}\right)
  =-\frac{1}{2\pi\sqrt{-1}}\rho\left(\frac{d}{d\tau}\right).
\]
以下我们需要考虑椭圆曲线族的退化。这时我们将使用参数
\[
  q=\exp(2\pi\sqrt{-1}\tau).
\]
于是我们有
\[
\begin{aligned}
  q\frac{d}{dq} &= \frac{1}{2\pi\sqrt{-1}}\frac{d}{d\tau}\\
  \theta\left(u^{-1}\frac{d}{du}\right) &= -\theta\left(q\frac{d}{dq}\right).
\end{aligned}
\]
这个公式将在后面用到。
\end{example}

% [Tutor Note: 这两行是后面（§1.3 附近及第 3 章）讨论退化椭圆曲线时反复使用的
%   换算：$q=e^{2\pi\sqrt{-1}\tau}$ 是模参数，$\frac{d}{d\tau}=2\pi\sqrt{-1}q\frac{d}{dq}$，
%   于是 $\rho(\frac{d}{d\tau})=2\pi\sqrt{-1}\rho(q\frac{d}{dq})$。原文 p.8
%  （另）三个概念须分清：完备（complete）= 任何别的族的纤维都能被拉回到本族
%   （存在 $g$ 即可）；半通用（versal）= 完备且 $g$ 的一阶微分 $(dg)_s$ 唯一，
%   即切空间层面唯一；通用（universal）= $g$ 本身唯一。原文"versal"常译"通用"，
%   这里按 STYLE 区分为半通用/通用。原文 p.8]
% [Tutor Note: (1.7) 图的可交换性是说：先沿纤维内做 $f$ 再投影，等于先投影到 $U$
%   再取 $g$；即 $g$ 就是 $f$ 在底空间上诱导的映射。原文这里只画了方图，
%   $\varpi^{-1}(U)\to U$ 与 $\mathcal{C}\to\cW$ 分别是两族的投影。原文 p.8
%  （另）定理 1.6 的 $g\ge 2$ 不可去掉：$g=1$ 时椭圆曲线有平移自同构、
%   $g=0$ 时有 $PGL(2,\C)$，正是"半通用而不通用"的障碍（该点处的自同构群非平凡，
%   见定义 1.5 后的那句话）。原文 p.8–p.9]

\noindent\textbf{1.1.4 Kuranishi 族.}\quad 我们引入下面的定义，以帮助我们理解小平–Spencer 映射的另一个方面。

\begin{definition}[1.5]
若紧黎曼曲面的复解析族 $\pi:\mathcal{C}\to\cW$ 与点 $0\in\cW$（$C=\pi^{-1}(0)$）满足下述条件 $(*)$，则称族 $\pi:\mathcal{C}\to\cW$ 在点 $0$ 处是\emph{完备的}（complete）。

$(*)$ 若给定紧黎曼曲面的复解析族 $\varpi:\mathcal{M}\to\mathcal{N}$、点 $s\in\mathcal{N}$ 以及解析同构 $f_0:M=\varpi^{-1}(s)\to C$，则存在从点 $s$ 的邻域 $U$ 到 $\cW$ 的全纯映射 $g$，以及全纯映射 $f:\varpi^{-1}(U)\to\mathcal{C}$，使得它们满足以下条件：
\[
g(s)=0,\qquad f|_{\varpi^{-1}(s)}=f_0,
\]
并且使得对任意点 $u\in U$，$f|_{\varpi^{-1}(u)}$ 都是从 $\varpi^{-1}(u)$ 到 $\pi^{-1}(g(u))$ 的解析同构，且有交换图
\begin{equation}
\begin{array}{ccc}
\varpi^{-1}(U) & \xrightarrow{\ f\ } & \mathcal{C}\\[1mm]
\Big\downarrow & & \Big\downarrow\\[1mm]
U & \xrightarrow{\ g\ } & \cW
\end{array} \tag{1.7}
\end{equation}

此时，族 $\varpi^{-1}(U)\to U$ 复解析同构于族 $\pi:\mathcal{C}\to\cW$ 沿全纯映射 $g$ 的拉回。此外，若对所有满足上述条件的 $g$，$(dg)_s$ 被唯一确定，则称族 $\pi:\mathcal{C}\to\cW$ 在点 $0$ 处是\emph{半通用的}（versal）。

此外，若 $g$ 被唯一确定，则称族 $\pi:\mathcal{C}\to\cW$ 在点 $0$ 处是\emph{通用的}（universal）。
\end{definition}

若解析族 $\pi:\mathcal{C}\to\cW$ 在点 $0$ 处是半通用但不通用的，则紧黎曼曲面 $C=\pi^{-1}(0)$ 的解析自同构群给出成为通用的障碍。下述定理在应用中很重要。

\begin{theorem}[1.6]
若亏格 $g\ge2$ 的紧黎曼曲面的复解析族 $\pi:\mathcal{C}\to\cW$ 在点 $0\in\cW$ 的某邻域内的每一点处都是完备的，且在 $0$ 处是半通用的，则族 $\pi$ 在点 $0$ 处是通用的。
\end{theorem}

% [原文 p.9]
小平–Spencer 映射与完备族或通用族之间的关系由下述定理给出。

\begin{theorem}[1.7]
若复解析族 $\pi:\mathcal{C}\to\cW$ 的小平–Spencer 映射在点 $w\in\cW$ 处是满射，则该族在点 $w$ 处是完备的。
\end{theorem}

\begin{theorem}[1.8]
复解析族 $\pi:\mathcal{C}\to\cW$ 在点 $0\in\cW$ 处是半通用的，当且仅当该族在点 $0$ 的某邻域内的每一点处都是完备的，且点 $0$ 处的小平–Spencer 映射 $\rho_0:T_0\cW\to H^1(\pi^{-1}(0),\Theta)$ 是同构。
\end{theorem}

\noindent\textbf{1.1.5 周期映射.}\quad 设 $R$ 是亏格为 $g$ 的紧黎曼曲面，$\{\alpha_1,\dots,\alpha_g,\beta_1,\dots,\beta_g\}$ 是 $H_1(R,\Z)$ 的辛基。黎曼曲面 $R$ 的周期矩阵（period matrix）$\tau$ 由
\begin{equation}
  \tau:=\left(\int_{\beta_i}\omega_j\right)\cdot\left(\int_{\alpha_i}\omega_j\right)^{-1} \tag{1.8}
\end{equation}
给出，其中 $\{\omega_1,\dots,\omega_g\}$ 是 $R$ 上全纯 1-形式的一组基。注意周期矩阵 $\tau$ 不依赖于全纯 1-形式的基的选取，但依赖于辛基的选取。设 $\{\tilde\alpha_1,\dots,\tilde\alpha_g,\tilde\beta_1,\dots,\tilde\beta_g\}$ 是 $H_1(R,\Z)$ 的另一组辛基。那么我们有
\[
  \begin{pmatrix}
    \tilde\beta_1\\ \vdots\\ \tilde\beta_g\\ \tilde\alpha_1\\ \vdots\\ \tilde\alpha_g
  \end{pmatrix}
  =
  \begin{pmatrix}
    A & B\\ C & D
  \end{pmatrix}
  \begin{pmatrix}
    \beta_1\\ \vdots\\ \beta_g\\ \alpha_1\\ \vdots\\ \alpha_g
  \end{pmatrix},
\]
其中矩阵
\[
  \begin{pmatrix}
    A & B\\ C & D
  \end{pmatrix}
\]
属于由
\[
  \operatorname{Sp}(g,\Z):=\left\{X\in M_{2g}(\Z)\ \middle|\ {}^tX
  \begin{pmatrix}
    0 & -I_g\\ I_g & 0
  \end{pmatrix}
  X=
  \begin{pmatrix}
    0 & -I_g\\ I_g & 0
  \end{pmatrix}
  \right\}
\]
定义的辛整数矩阵群 $\operatorname{Sp}(g,\Z)$。设 $\tilde\tau$ 是由辛基 $\{\tilde\alpha_1,\dots,\tilde\alpha_g,\tilde\beta_1,\dots,\tilde\beta_g\}$ 确定的黎曼曲面 $R$ 的周期矩阵。那么我们有
\[
  \tilde\tau=(A\tau+B)(C\tau+D)^{-1}.
\]
$R$ 上全纯 1-形式的一组基 $\{\omega_1,\dots,\omega_g\}$ 称为关于 $H_1(R,\Z)$ 的辛基 $\{\alpha_1,\dots,\alpha_g,\beta_1,\dots,\beta_g\}$ 的标准化基（normalized basis），若
\begin{equation}
  \int_{\alpha_i}\omega_j=\delta_{ij},\qquad \int_{\beta_i}\omega_j=\tau_{ij},\qquad \tau=(\tau_{ij}). \tag{1.9}
\end{equation}
设 $\pi:\mathcal{R}\to\mathbf{T}$ 是亏格为 $g$ 的紧黎曼曲面在 Teichmüller 空间 $\mathbf{T}$ 上的通用族。对每一点 $t\in\mathbf{T}$，置 $R_t=\pi^{-1}(t)$。则存在典范同构
\begin{equation}
  H^1(R_t,\Theta_{R_t})\simeq T_t\mathbf{T}. \tag{1.10}
\end{equation}

% [Tutor Note: (1.8) 中两个积分矩阵的形状：$\left(\int_{\beta_i}\omega_j\right)$ 的第 $i$
%   行第 $j$ 列是 $\int_{\beta_i}\omega_j$（$i$ 为周期、$j$ 为形式），右乘
%   $\left(\int_{\alpha_i}\omega_j\right)^{-1}$ 后得到 $g\times g$ 矩阵 $\tau$；这正是
%   周期矩阵公式 $\tau=\Pi_B\Pi_A^{-1}$。原文 p.9
%  （另）术语首次出现：period matrix = 周期矩阵；normalized basis = 标准化基；
%   Teichmüller space = Teichmüller 空间；symplectic integral matrices 按原文符号
%   $\operatorname{Sp}(g,\Z)$ 保留。$\mathbf{T}$ 是粗排的粗体，不是层记号。原文 p.9]
% [原文 p.10]
它不是别的，正是由小平–Spencer 映射诱导的同构。取其对偶，得到典范同构
\begin{equation}
  T_t\mathbf{T}^*\simeq H^0(R_t,\omega_{R_t}^{\otimes 2}). \tag{1.11}
\end{equation}
以下我们常常把 (1.11) 的两边等同起来。由于 $\mathbf{T}$ 是单连通的，层 $R_1\pi_*\Z$ 在 $\mathbf{T}$ 上是平凡的，于是我们得到层 $R_1\pi_*\Z$ 的整体截面
\[
  \{\hat\alpha_1,\dots,\hat\alpha_g,\hat\beta_1,\dots,\hat\beta_g\},
\]
使得对每一点 $t\in\mathbf{T}$，
\[
  \{\hat\alpha_1(t),\dots,\hat\alpha_g(t),\hat\beta_1(t),\dots,\hat\beta_g(t)\}
\]
是 $H_1(R_t,\Z)$ 的辛基。层 $\pi_*\omega_{\mathcal{R}/\mathbf{T}}$ 是秩为 $g$ 的自由 $\cO_{\mathbf{T}}$-模。设 $\{\hat\omega_1,\dots,\hat\omega_g\}$ 是 $\pi_*\omega_{\mathcal{R}/\mathbf{T}}$ 的 $\cO_{\mathbf{T}}$-自由基。那么矩阵
\begin{equation}
  \Pi(t):=\left(\int_{\hat\beta_i(t)}\hat\omega_j(t)\right)\cdot
  \left(\int_{\hat\alpha_i(t)}\hat\omega_j(t)\right)^{-1} \tag{1.12}
\end{equation}
定义全纯映射
\begin{equation}
  \Pi:\mathbf{T}\longrightarrow\mathscr{S}_g \tag{1.13}
\end{equation}
（从 Teichmüller 空间到 $g$ 次 Siegel 上半空间）。此外，存在群同态
\begin{equation}
  \Phi:\operatorname{Mod}_g\longrightarrow \operatorname{Sp}(g,\Z) \tag{1.14}
\end{equation}
（从亏格 $g$ 的映射类群 $\operatorname{Mod}_g$ 到辛群），使得
\begin{equation}
  \Pi(\gamma t)=(A_\gamma\Pi(t)+B_\gamma)(C_\gamma\Pi(t)+D_\gamma)^{-1}, \tag{1.15}
\end{equation}
其中
\begin{equation}
  \Phi(\gamma)=
  \begin{pmatrix}
    A_\gamma & B_\gamma\\ B_\gamma & D_\gamma
  \end{pmatrix}. \tag{1.16}
\end{equation}

% [Tutor Note: 原文用花体/德文尖角体写 $g$ 次 Siegel 上半空间，本译注本无 eufrak 字体，
%   改用 \texttt{\textbackslash mathscr\{S\}}（mathrsfs 已载入）；记号意义不变：$\mathscr{S}_g$
%   $=\{\tau\in M_g(\C)\mid \tau=\tau^t,\ \operatorname{Im}\tau>0\}$，即周期矩阵所在的空间。
%   原文 p.10
%  （另）笔误：(1.16) 的第二行第一项印作 $B_\gamma$，应为 $C_\gamma$——矩阵
%   $\Phi(\gamma)$ 的分块按 (1.15) 的定义是 $\begin{psmallmatrix}A_\gamma&B_\gamma\\
%   $(A_\gamma,B_\gamma;C_\gamma,D_\gamma)$ 属于 $\operatorname{Sp}(g,\Z)$；若第一列为
%   $(A_\gamma,B_\gamma)^t$ 就会与 (1.15) 中 $C_\gamma$ 的用法冲突。照原文排，不改。原文 p.10]

\noindent\textbf{1.1.6 黎曼曲面的雅可比簇.}\quad 对紧黎曼曲面 $R$，设 $\{\alpha_1,\dots,\alpha_g,\beta_1,\dots,\beta_g\}$ 是辛基，$\{\omega_1,\dots,\omega_g\}$ 是全纯 1-形式的标准化基。以下
\[
  (R,\{\alpha_1,\dots,\alpha_g,\beta_1,\dots,\beta_g\})
\]
称为弱标记黎曼曲面（weakly marked Riemann surface）。设 $\tau\in\mathscr{S}_g$ 是相应的周期矩阵。我们用 $J(R)$ 表示以 $(I_g,\tau)$ 为周期矩阵的 $g$ 维复环面。那么我们可以定义全纯映射 $j:R\to J(R)$：
\begin{equation}
  j(Q):=\left(\int_{P_0}^{Q}\omega_1,\dots,\int_{P_0}^{Q}\omega_g\right), \tag{1.17}
\end{equation}
其中基点 $P_0\in R$ 一次取定后不再变动。若改变基点，映射 $j$ 只差复环面 $J(R)$ 的一个平移。映射 $j$ 称为 Albanese 映射，而复环面 $J(R)$ 通常称为黎曼曲面 $R$ 的雅可比簇（Jacobian variety）。

雅可比簇 $J(R)$ 也可以用内蕴的方式定义为
\begin{equation}
  J(R)=H^0(R,\omega_R)^*/H_1(R,\Z). \tag{1.18}
\end{equation}
确切地说，这是黎曼曲面 $R$ 的 Albanese 簇的定义。由于我们有
\[
  \int_{\alpha_i}\omega_j=\delta_{ij},
\]

% [原文 p.11]
故对 $H^0(R,\omega_R)$ 的标准化基 $\{\omega_1,\dots,\omega_g\}$，我们可以把 $\{\alpha_1,\dots,\alpha_g\}$ 看作 $H^0(R,\omega_R)^*$ 的相应对偶基。于是等式
\[
  \int_{\beta_i}\omega_j=\tau_{ij}
\]
蕴含
\[
  \beta_i=\sum_{j=1}^{g}\tau_{ij}\alpha_j .
\]
因此，我们可以把 $J(R)$ 表示成上面定义的复环面 $\C^g/(I_g,\tau)$。

黎曼曲面 $R$ 的 Picard 簇（Picard variety）$Pic^0(R)$ 定义为
\begin{equation}
  Pic^0(R):=H^1(R,\cO_R)/H^1(R,\Z). \tag{1.19}
\end{equation}
再次取 $H^0(R,\omega_R)$ 的标准化基 $\{\omega_1,\dots,\omega_g\}$。那么 $\{\ol\omega_1,\dots,\ol\omega_g\}$ 是 $H^1(R,\cO_R)$ 的一组基。由 Hodge 分解定理，我们有同构
\[
  H^1(R,\C)\simeq H^1(R,\cO_R)\oplus H^0(R,\omega_R).
\]
\begin{sloppypar}设 $\{\alpha_1^*,\dots,\alpha_g^*,\beta_1^*,\dots,\beta_g^*\}$ 是上同调群 $H^1(R,\Z)$ 的对偶于 $H_1(R,\Z)$ 的基 $\{\alpha_1,\dots,\alpha_g,\beta_1,\dots,\beta_g\}$ 的基。由于我们可以把 $H^1(R,\Z)$ 看作 $H^1(R,\C)$ 的格，由上述同构我们可以把 $\alpha_i^*$、$\beta_i^*$ 表示为
\end{sloppypar}
\[
\begin{aligned}
  \alpha_i^* &= \sum_{j=1}^{g}a_{ij}\ol\omega_j+\sum_{j=1}^{g}\ol a_{ij}\omega_j,\\
  \beta_i^* &= \sum_{j=1}^{g}b_{ij}\ol\omega_j+\sum_{j=1}^{g}\ol b_{ij}\omega_j .
\end{aligned}
\]
置
\[
  A=(a_{ij}),\qquad B=(b_{ij}).
\]
由等式
\[
\begin{aligned}
  \delta_{ik}=\langle\alpha_k,\alpha_i^*\rangle
  &= \sum_{j=1}^{g}a_{ij}\int_{\alpha_k}\ol\omega_j+\sum_{j=1}^{g}\ol a_{ij}\int_{\alpha_k}\omega_j\\
  &= a_{ik}+\ol a_{ik},\\[1mm]
  0=\langle\beta_k,\alpha_i^*\rangle
  &= \sum_{j=1}^{g}a_{ij}\int_{\beta_k}\ol\omega_j+\sum_{j=1}^{g}\ol a_{ij}\int_{\beta_k}\omega_j\\
  &= \sum_{j=1}^{g}a_{ij}\ol\tau_{jk}+\sum_{j=1}^{g}\ol a_{ij}\tau_{jk}
\end{aligned}
\]
以及关于 $\beta_i^*$ 的类似等式，我们有
\[
  \begin{pmatrix} A & \ol A\\ B & \ol B\end{pmatrix}
  \begin{pmatrix} I_g & \ol\tau\\ I_g & \tau\end{pmatrix}=I_{2g}.
\]
置
\begin{equation}
  Y=(-2\pi\sqrt{-1}\operatorname{Im}\tau)^{-1}. \tag{1.20}
\end{equation}
注意 $Y$ 是对称矩阵。我们有
\[
  \begin{pmatrix} I_g & \ol\tau\\ I_g & \tau\end{pmatrix}^{-1}
  =\begin{pmatrix} -\tau Y & -\ol\tau Y\\ Y & Y\end{pmatrix}.
\]

% [Tutor Note: 这一步的核对方法：把 $\alpha_i^*$ 按 Hodge 分解写成 $\sum a_{ij}\ol\omega_j+\sum\ol a_{ij}\omega_j$，
%  再用对偶配对 $\langle\alpha_k,\omega_j\rangle=\int_{\alpha_k}\omega_j=\delta_{kj}$、
%  $\langle\beta_k,\omega_j\rangle=\tau_{kj}$；注意 $\langle\alpha_k,\ol\omega_j\rangle=\ol{\langle\alpha_k,\omega_j\rangle}=\delta_{kj}$，
%  所以 $\delta_{ik}=a_{ik}+\ol a_{ik}$（这正是 $\alpha_i^*$ 的实性条件），而不是 $a_{ik}$ 单独。
%  原文 p.11]
% [原文 p.12]
于是，由 $H^1(R,\Z)$ 在 $H^1(R,\cO_R)$ 中诱导的格由矩阵
\[
  \begin{pmatrix} -\tau Y\\ Y\end{pmatrix}
\]
的行给出，其中我们以 $\{\ol\omega_1,\dots,\ol\omega_g\}$ 作为 $H^1(R,\cO_R)$ 的基。因此，我们可以把 $Pic^0(R)$ 表示成复环面 $\C^g/(Y\tau,Y)$。于是存在同构
\[
  \lambda:J(R)\longrightarrow Pic^0(R),
\]
它由
\begin{equation}
  \lambda\begin{bmatrix} z_1\\ \vdots\\ z_g\end{bmatrix}
  =\begin{bmatrix} Y\begin{pmatrix} z_1\\ \vdots\\ z_g\end{pmatrix}\end{bmatrix} \tag{1.21}
\end{equation}
给出。注意同构 $\lambda$ 是由 $\C$-线性同构
\begin{equation}
\begin{aligned}
  H^0(R,\omega_R)^* &\longrightarrow H^1(R,\cO_R),\\
  \sum a_i\alpha_i &\longmapsto \sum_{ij}(y_{ij}a_j)\ol\omega_i
\end{aligned} \tag{1.22}
\end{equation}
诱导的，其中 $(y_{ij})=Y$。

紧黎曼曲面的 Albanese 簇与 Picard 簇互为对偶，且它们都是主极化的阿贝尔簇。上述同构 $\lambda:J(R)\to Pic^0(R)$ 是典范的。雅可比簇上全纯向量场的层 $\Theta_{J(R)}$ 是平凡的，且我们有典范同构
\[
  \Theta_{J(R)}\simeq T_0J(R)\otimes\cO_{J(R)}.
\]
此外，Albanese 映射 $j:R\to J(R)$ 通过
\[
  d\ol z_j\longmapsto \ol\omega_j
\]
诱导同构
\[
  H^1(J(R),\cO)\simeq H^1(R,\cO).
\]
因此，我们有典范同构
\[
  \iota:H^1(J(R),\Theta_{J(R)})\simeq T_0J(R)\otimes H^1(R,\cO_R),
\]
其中 $T_0J(R)$ 是 $J(R)$ 在原点处的切空间。由上述同构 $\lambda$，我们可以把 $T_0J(R)$ 与 $H^1(R,\cO_R)$ 等同起来。于是我们可以把上式改写为
\begin{equation}
  \iota:H^1(J(R),\Theta_{J(R)})\simeq H^1(R,\cO_R)\otimes H^1(R,\cO_R). \tag{1.23}
\end{equation}
由于 $T_0J(R)$ 被 (1.19) 与 $H^0(R,\omega_R)^*$ 等同，由 (1.8)) 映射 $\iota$ 可显式地写成
\begin{equation}
  \iota\left(\frac{\partial}{\partial z_i}d\ol z_j\right)
  =\sum_{k=1}^{g}y_{ik}\ol\omega_k\otimes\ol\omega_j, \tag{1.24}
\end{equation}
其中矩阵 $(y_{ij})=Y$ 由 (1.20) 定义。置
\[
  \mathscr{T}:=\left\{T=(t_{ij})\ \middle|\ T\ \text{是}\ g\times g\ \text{复矩阵，且}\
  \det\begin{pmatrix} I_g & T\\ I_g & \ol T\end{pmatrix}\ne 0\right\},
\]
并让 $\Z^{2g}$ 作用在 $\mathscr{T}\times\C^g$ 上：
\[
  n:(T,z)\longmapsto (T,z+n_1+n_2T),\qquad n=(n_1,n_2),\quad n_i\in\Z^g .
\]
商空间 $\pi:\mathcal{T}=\mathscr{T}\times\C^g/\Z^{2g}\to\mathscr{T}$ 是复环面的通用族。设 $\mu:\mathcal{A}\to\mathscr{S}_g$ 是 Siegel 上半空间上的阿贝尔簇族，使得纤维 $A_\tau=\mu^{-1}(\tau)$ 是以 $(I_g,\tau)$ 为周期矩阵的复环面。则

% [原文 p.13]
存在一个自然嵌入 $i:\mathscr{S}_g\hookrightarrow\mathscr{T}$，且诱导族 $i^*\mathcal{T}\to\mathscr{S}_g$ 同构于族 $\mu:\mathcal{A}\to\mathscr{S}_g$。我们把映射
\[
  i\circ\Pi:\mathbf{T}\longrightarrow\mathscr{T}
\]
也称为周期映射。在与具有辛基 $\{\alpha_1,\dots,\alpha_g,\beta_1,\dots,\beta_g\}$ 的黎曼曲面 $R$ 相对应的点 $\tau\in\mathscr{S}_g\subset\mathscr{T}$ 处，借小平–Spencer 映射我们可以把 $\mathscr{T}$ 在 $\tau$ 处的切空间 $T_\tau\mathscr{T}$ 与 $H^1(J(R),\Theta_{J(R)})$ 等同起来。因此，周期映射 $i\circ\Pi$ 诱导出自然映射
\[
  H^1(R,\Theta_R)\longrightarrow H^1(J(R),\Theta_{J(R)})=H^1(R,\cO_R)\otimes H^1(R,\cO_R),
\]
其中我们用了同构 $\iota$ (1.22)。该映射的对偶由自然的乘法映射
\begin{equation}
\begin{aligned}
  H^0(R,\omega_R)\otimes H^0(R,\omega_R)&\longrightarrow H^0(R,\omega_R^{\otimes2})\\
  \nu\otimes\mu&\longmapsto \nu\mu
\end{aligned} \tag{1.25}
\end{equation}
给出。

下面我们计算点 $\tau\in\mathscr{S}_g\subset\mathscr{T}$ 处的切向量 $\frac{\partial}{\partial t_{ij}}$ 的小平–Spencer 类。设 $\mathcal{U}$ 是点 $\tau_0$ 的开邻域，$\{\cU_\lambda\}_{\lambda\in\Lambda}$ 是 $\mathcal{U}$ 的小开覆盖。我们把局部坐标 $(t_{ij},z_\lambda^1,\dots,z_\lambda^g)$（$1\le i\le j\le g$）取成使得 $(z_\lambda^1,\dots,z_\lambda^g)$ 由向量空间 $\C^g$ 的整体仿射坐标 $(z_1,\dots,z_g)$ 导出。于是，若 $\cU_\lambda\cap\cU_\mu\ne\emptyset$，则我们有
\begin{equation}
  z_\lambda^i-z_\mu^i=m_{\lambda\mu}^i+\sum_{k=1}^{g}n_{\lambda\mu}^k t_{ik} \tag{1.26}
\end{equation}
其中 $m_{\lambda\mu}^i,n_{\lambda\mu}^i\in\R^g$，$i=1,2,\dots,g$。因此，小平–Spencer 类 $\rho\left(\frac{\partial}{\partial t_{ij}}\right)$ 由上链 $\theta_{\lambda\mu}$ 给出，它定义为
\begin{equation}
  \theta_{\lambda\mu}=n_{\lambda\mu}^j\frac{\partial}{\partial z_\lambda^i}. \tag{1.27}
\end{equation}
现在置
\[
  \begin{pmatrix}
    \psi_\lambda^1\\
    \vdots\\
    \psi_\lambda^g
  \end{pmatrix}
  =Y\begin{pmatrix}
    \ol z_\lambda^1-z_\lambda^1\\
    \vdots\\
    \ol z_\lambda^g-z_\lambda^g
  \end{pmatrix}.
\]
其中 $Y=(-2\sqrt{-1}\operatorname{Im}(t_{ij}))^{-1}$。那么我们有
\[
  \psi_\mu^j\frac{\partial}{\partial z_\mu^j}-\psi_\lambda^j\frac{\partial}{\partial z_\lambda^j}=\theta_{\lambda\mu}.
\]
因此
\[
  \ol\partial\left(\psi_\mu^j\frac{\partial}{\partial z_\mu^j}\right)
  =\ol\partial\left(\psi_\lambda^j\frac{\partial}{\partial z_\lambda^j}\right)
\]
是整体形式，它表示为
\begin{equation}
  \sum_{k=1}^{g}y_{jk}\frac{\partial}{\partial z_i}d\ol z_k. \tag{1.28}
\end{equation}

% [Tutor Note: 笔误：本段"由 (1.8))"多了一个右括号，且按上下文 $T_0J(R)\simeq H^0(R,\omega_R)^*$
%   用的是 (1.18)（$J(R)=H^0(R,\omega_R)^*/H_1(R,\Z)$），(1.8) 是周期矩阵的定义。
%   照原文排，不改；读者按 (1.18) 理解即可。原文 p.12
%  （另）$\mathscr{T}$ 中 $\det\begin{psmallmatrix}I_g&T\\I_g&\ol T\end{psmallmatrix}\ne0$
%   等价于 $\tau=T$ 时 $\operatorname{Im}\tau>0$ 的推广：当 $T=\tau$ 对称时该行列式为
%   $\det(\ol\tau-\tau)\ne0$，即 $\operatorname{Im}\tau$ 非退化。$\Z^{2g}$ 作用正是
%   $(I_g,T)$ 的整格平移，故 $\mathcal{T}$ 是"所有复环面+极化"的模空间。原文 p.12
%  （另）笔误：(1.27) 中 $n_{\lambda\mu}^j$ 的上标 $j$ 与 $\partial/\partial z_\lambda^i$ 的 $i$
%   不一致——$\rho(\partial/\partial t_{ij})$ 应配 $n_{\lambda\mu}^i\partial/\partial z_\lambda^i$
%   （对照 (1.26)：$z_\lambda^i-z_\mu^i$ 展开出的系数是 $n_{\lambda\mu}^k$ 里 $k$ 与 $t_{ik}$
%   的第二个指标缩并，故留下的是 $t_{ij}$ 的 $j$ 对应 $n^j$ 而 $z^i$ 与 $i$ 对应）。
%   照原文排，不改。原文 p.13
%  （另）(1.26) 中 $n_{\lambda\mu}^k$ 印作上标 $k$（而非 $ik$），与 (1.27) 的写法配套；
%   这里 $t_{ik}$ 是周期矩阵的对称坐标（$i\le k$），$z_\lambda^i$ 是纤维上的坐标。原文 p.13
%  （另）(1.28) 里 $\partial/\partial z_i$ 的下标没有 $\lambda$（与前一段的
%   $\partial/\partial z_\lambda^i$ 不同），疑为原文排印简化；照排不改。原文 p.13]
% [原文 p.14]
Dolbeault 上同调类 (1.0-20) 就是小平–Spencer 类 $\rho\left(\frac{\partial}{\partial t_{ij}}\right)$，并且由 (1.22) 我们有
\[
  \iota\left(\rho\left(\frac{\partial}{\partial t_{ij}}\right)\right)
  =\sum_{\ell=1}^{g}\sum_{k=1}^{g}\{(y_{i\ell}\ol\omega_\ell)\otimes(y_{jk}\ol\omega_k)\}.
\]
于是，由 (1.22)，这个类在 $H^0(R,\omega_R)\otimes H^0(R,\omega_R)$ 中的对偶是 $\omega_i\otimes\omega_j$。因此，由 (1.26) 我们断言
\begin{equation}
  \Pi^*(d\tau_{ij})=\Pi^*\circ i^*(dt_{ij})=\omega_i\omega_j, \tag{1.29}
\end{equation}
其中我们用了等同 (1.11)。这样我们就证明了下列命题。

\begin{proposition}[1.9]
在上述假定与记号下，我们有
\begin{equation}
\begin{aligned}
  &\Pi^*\left(\sum_{i<j}\frac{\partial}{\partial\tau_{ij}}\left(\log\det(C_\gamma\tau+D_\gamma)\right)d\tau_{ij}\right)\\
  &\qquad =\sum_{i<j}\omega_i\omega_j\Pi^*\left(\frac{\partial}{\partial\tau_{ij}}\log\det(C_\gamma\tau+D_\gamma)\right).
\end{aligned} \tag{1.30}
\end{equation}
\end{proposition}

\noindent\textbf{1.1.7 对称双微分.}\quad 我们考虑紧黎曼曲面上的对称双微分形式。首先我们回顾 theta 函数的基本事实。对点 $\tau\in\mathscr{S}_g$，theta 函数 $\theta(\tau,z)$ 定义为
\[
  \theta(\tau,z):=\sum_{m\in\Z^g}\exp\left\{2\pi\sqrt{-1}\left(\frac{1}{2}{}^tm\tau m+{}^tmz\right)\right\}.
\]
theta 函数 $\theta(\tau,z)$ 可以看作定义了阿贝尔簇 $A_\tau=\C^g/(I_g,\tau)$ 的主极化的线丛的全纯截面。对 $e\in\C^g$，我们用 $[e]$ 表示 $A_\tau$ 中对应的点。theta 函数 $\theta(\tau,z)$ 的零点定义 $A_\tau$ 上的一个除子，称为 \emph{theta 除子}（theta divisor），记作 $\Theta$。

具有特征 $e=\begin{bmatrix}\delta\\ \epsilon\end{bmatrix}$（$\delta,\epsilon\in\R^g$）的 theta 函数 $\theta[e](\tau,z)$ 定义为
\[
  \theta[e](\tau,z):=\sum_{m\in\Z^g}\exp 2\pi\sqrt{-1}\left\{\frac{1}{2}{}^t(m+\delta)\tau(m+\delta)+{}^t(m+\delta)(z+\epsilon)\right\}.
\]

设 $R$ 是亏格为 $g$ 的紧黎曼曲面。固定 $H_1(R,\Z)$ 的一组辛基，并设 $\tau\in\mathscr{S}_g$ 是对应的周期矩阵。那么阿贝尔簇 $A_\tau$ 就是 $R$ 的 Jacobi 簇 $J(R)$。设 $Div^d(R)$ 是紧黎曼曲面 $R$ 上 $d$ 次除子的群。我们可以把映射 $j:R\to J(R)$ 扩张为映射 $j:Div^d(R)\to J(R)$，为此置
\[
  j\left(\sum_{i=1}^{m+d}P_i-\sum_{k=1}^m Q_k\right)
  :=\sum_{i=1}^{m+d}j(P_i)-\sum_{k=1}^m j(Q_k).
\]
注意对两个除子 $D,D'\in Div^d(R)$，$j(D)=j(D')$ 当且仅当 $D$ 与 $D'$ 线性等价，即存在 $R$ 上的亚纯函数 $f$ 使得 $D=D'+(f)$。设 $Pic^d(R)$ 是紧黎曼曲面 $R$ 上 $d$ 次线丛的同构类之集。那么由 Abel 定理，映射 $j:Div^d(R)\to J(R)$ 穿过 $j:Pic^d(R)\to J(R)$，其中我们把 $Pic^d(R)$ 与商 $Div^d(R)/\sim$（按线性等价关系）等同起来。此外，对 $R$ 上任一除子 $D$，我们用 $[D]$ 表示 $R$ 上对应的线丛。下面我们总是通过典范同构 $\lambda$ 把 $Pic^0(R)$ 与 $J(R)$ 等同起来。下述由 Riemann 得到的定理是与黎曼曲面相联系的 theta 函数理论中最基本的定理。

% [Tutor Note: 笔误：本页开头 "The Dolbeault cohomology class (1.0-20)" 的编号
%   显然是排印错误，按上下文应指 (1.28)（该式给出 Dolbeault 形式
%   $\sum_{k=1}^g y_{jk}\frac{\partial}{\partial z_i}d\ol z_k$）。照原文排，不改。原文 p.14
%  （另）术语首次出现：theta divisor = theta 除子；characteristic = 特征；
%   $Div^d(R)$、$Pic^d(R)$ 按原文保留斜体，不排成算子名。原文 p.14]

% [原文 p.15]
\begin{theorem}[1.10]
存在除子 $\Delta\in Div^{g-1}(R)$，使得
\begin{equation}
  [2\Delta]=\omega_R, \tag{1.31}
\end{equation}
且对任意点 $P\in R$ 与 $e\in\C^g$，我们有：
\begin{enumerate}[label=(\arabic*)]
\item 若 $\theta(\tau,e)\ne0$，则 $R$ 的由方程（变元为 $Q$）
\[
  \theta(\tau,j(Q)-j(P)-e)=0
\]
定义的除子 $D$ 是 $g$ 次有效除子，且
\[
  H^1(R,\cO_R(D))=0,\qquad [e]=j(D-P-\Delta).
\]
\item 若 $\theta(\tau,e)=0$，则对某个 $g-1$ 次有效除子 $E$ 我们有
\[
  [e]=j(E-\Delta).
\]
\end{enumerate}
此外，$\dim_\C H^1(R,\cO_R(E))$ 是 theta 除子 $\Theta$ 在点 $[e]$ 处的重数，也是使得
\[
  \theta(B-A-e)\equiv0
\]
对一切次数小于或等于 $d-1$ 的有效除子 $A,B$ 成立的最小整数 $d$。
\end{theorem}

置
\begin{equation}
  L_0=[\Delta]. \tag{1.32}
\end{equation}
于是我们有
\begin{equation}
  L_0^{\otimes2}=\omega_R. \tag{1.33}
\end{equation}
对半周期（即二阶挠点）$\alpha\in J(C)$，也就是说 $2\alpha=0$ 于 $J(R)$，设 $L_\alpha$ 是 $R$ 上的线丛，使得 $L_\alpha\otimes L_0^{-1}$ 是半周期 $\alpha$。于是特别地
\begin{equation}
  L_\alpha^{\otimes2}=\omega_R. \tag{1.34}
\end{equation}

\begin{corollary}[1.11]
我们有
\begin{equation}
  h^0(R,L_\alpha)=h^1(R,L_\alpha)=\operatorname{mult}_\alpha\Theta, \tag{1.35}
\end{equation}
且 $h^0(R,L_\alpha)$ 在 $\alpha$ 为偶时是偶数、在 $\alpha$ 为奇时是奇数，其中半周期 $\alpha$ 称为奇（相应地，偶）的，若我们有
\[
  \alpha=\frac{1}{2}\begin{bmatrix}\delta\\ \epsilon\end{bmatrix}\begin{pmatrix}\tau\\ I_g\end{pmatrix},\qquad \delta,\epsilon\in\Z^g
\]
且 $\delta\cdot{}^t\epsilon$ 为偶（相应地，奇）。
\end{corollary}

\begin{sloppypar}定理 1.10 与推论 1.11 的证明我们请读者参阅 [Fay]。设 $\alpha$ 是非奇异的奇半周期，即 $\operatorname{mult}_\alpha\Theta=1$。那么由上述推论，线丛 $L_\alpha$ 有非零的整体截面 $h_\alpha$，满足
\end{sloppypar}
\begin{equation}
  h_\alpha^2=\sum_{i=1}^{g}\frac{\partial\theta[\alpha]}{\partial z_i}(\tau,0)\omega_i, \tag{1.36}
\end{equation}
其中 $\{\omega_1,\dots,\omega_g\}$ 是全纯 1-形式的标准化基，并且我们用了 (1.34)。

% [Tutor Note: 笔误：原文把 $J(R)$ 印成了 $J(C)$（"$\alpha\in J(C)$"），而本节
%   没有曲面 $C$，紧接着的 $2\alpha=0$ 也写在 $J(R)$ 里。照原文排，不改。原文 p.15
%  （另）(1.35) 中 $\operatorname{mult}_\alpha\Theta$ 是 theta 除子在点 $\alpha$ 处
%   的重数；$h^0=h^1$ 由 Serre 对偶（$\omega_R$ 与 $\cO_R$ 的配对）给出，
%   这才是"$\alpha$ 奇/偶 $\Longleftrightarrow$ $h^0$ 奇/偶"的来源。原文 p.15]

% [原文 p.16]
黎曼曲面 $R$ 的 \emph{素形式}（prime form）$E(P,Q)$ 定义为
\begin{equation}
  E(P,Q):=\frac{\theta[\alpha](\tau,j(P-Q))}{h_\alpha(P)h_\alpha(Q)} \tag{1.37}
\end{equation}
其中 $P,Q\in R$。注意素形式不依赖于 $\alpha$ 的选取，并且是 $R\times R$ 上线丛 $\pi_1^*L_0^{-1}\otimes\pi_2^*L_0^{-1}\otimes\delta^*(\Theta)$ 的全纯截面，其中 $\pi_i$ 是 $R\times R$ 到第 $i$ 个因子的投影，$\delta:R\times R\to J(R)$ 是把 $(P,Q)\in R\times R$ 送到 $j(Q-P)$ 的映射。此外，我们有

% [Tutor Note: (1.37) 的分母是 $h_\alpha(P)h_\alpha(Q)$，其中 $h_\alpha$ 是 (1.36)
%   给出的 $L_\alpha$ 的截面；由 (1.34) $L_\alpha^{\otimes2}=\omega_R$，$h_\alpha$
%   是 $\omega_R^{1/2}$ 型的半微分，故商 $\theta[\alpha]/(h_\alpha h_\alpha)$ 在
%   $P$、$Q$ 各自的坐标变换下不变——这才是"素形式不依赖于 $\alpha$ 的选取"的机制，
%   而不是省略号式的"可验证"。原文 p.16
%  （另）本小节后段的术语首次出现：biresidue = 双留数；projective connection =
%   射影联络；Schwarzian derivative = Schwarz 导数。原文 p.16–p.17]

\begin{equation}
  E(P,Q)=-E(Q,P), \tag{1.38}
\end{equation}
且 $E(P,Q)$ 沿 $R\times R$ 中的对角线一阶零点，在对角线之外不为零。设 $x$（相应地，$y$）是以 $P$（相应地，$Q$）为中心的 $R$ 的局部坐标，并滥用记号，让我们写
\[
  h_\alpha(P)=h_\alpha(x)\sqrt{dx},\qquad h_\alpha(Q)=h_\alpha(y)\sqrt{dy}.
\]
那么素形式 $E(P,Q)$ 表示为
\[
  E(P,Q)=E(x,y)(\sqrt{dy})^{-1}(\sqrt{dy})^{-1}.
\]
现在我们定义双微分 $\omega(x,y)dxdy$ 为
\begin{equation}
  \omega(x,y)dxdy:=\frac{\partial^2}{\partial x\partial y}\log E(x,y)dxdy. \tag{1.39}
\end{equation}
那么 $\omega(x,y)dxdy$ 是 $R\times R$ 上的亚纯双微分，在对角线之外全纯，并沿对角线有二阶极点。也就是说，$\omega(x,y)dxdy\in H^0(R\times R,p_1^*\omega_R\otimes p_2^*\omega_R(2\Delta))$。此外，由 (1.38) 我们有
\begin{equation}
  \omega(x,y)=\omega(y,x). \tag{1.40}
\end{equation}

\begin{proposition}[1.12]
对 $f\in\C^g$，若 $[f]\in J(R)$ 是 theta 除子 $\Theta$ 的非奇异点，则我们有
\[
  \omega(x,y)dxdy=\frac{\partial^2\log\theta}{\partial x\partial y}(\tau,j(x)-j(y)-f)dxdy.
\]
\end{proposition}
证明见 [Fay]。

对 $R\times R$ 上在对角线之外全纯的对称双微分 $\hat\omega(x,y)dxdy$，我们沿对角线总有展开
\[
  \hat\omega(x,y)dxdy=\left(\sum_{k=k_0}^{2}\frac{c_{-k}}{(x-y)^k}+\text{全纯部分}\right).
\]
系数 $c_{-2}$ 不依赖于局部坐标的选取，称为 \emph{双留数}（biresidue），记作 $\operatorname{Res}^2(\hat\omega)$。

\begin{corollary}[1.13]
存在对称双微分 $\omega$，它在对角线之外全纯，并沿对角线有二阶极点且 $\operatorname{Res}^2\omega=1$。
\end{corollary}

下面，对黎曼曲面 $R$ 的亚纯对称双微分 $\omega(x,y)dxdy$，我们总是指 $\omega(y,x)=\omega(x,y)$，$\omega$ 在 $R\times R$ 上对角线之外全纯，并沿对角线有二阶极点。

% [原文 p.17]
\begin{lemma}[1.14]
设 $\omega(x,y)dxdy$ 与 $\tilde\omega(x,y)dxdy$ 是 $R\times R$ 上在对角线之外全纯、沿对角线有二阶极点、且满足
\[
  \operatorname{Res}^2(\omega)=\operatorname{Res}^2(\tilde\omega)
\]
的对称双线性微分。那么差
\[
  \Omega(x,y)dxdy=\omega(x,y)dxdy-\tilde\omega(x,y)dxdy
\]
是 $R\times R$ 上的全纯对称双微分。
\end{lemma}

$R\times R$ 上沿对角线有二阶极点、在对角线之外全纯的对称双微分构成的空间记为
\[
  H^0(R\times R,p_1^*\omega_R\otimes p_2^*\omega_R(2\Delta))^{(\iota)},
\]
其中 $\iota$ 是 $R\times R$ 交换两个分支的自同构。注意
\[
  \omega_{R\times R}=p_1^*\omega_R\otimes p_2^*\omega_R .
\]
让我们对紧黎曼曲面 $R$ 上双留数为 1 的对称双微分 $\omega$，定义与它相伴的射影联络（projective connection）$S_\omega(z)dz^2$：
\begin{equation}
  S_\omega(z)dz^2=6\lim_{w\to z}\left(\omega(w,z)dwdz-\frac{dwdw}{(w-z)^2}\right). \tag{1.41}
\end{equation}
射影联络 $S_\omega(z)dz^2$ 确实依赖于局部坐标 $z$ 的选取。下述定理是众所周知的，可由直接计算得到。

\begin{theorem}[1.15]
\[
  S_\omega(z)dz^2=S_\omega(w)dw^2-\{z;w\}dw^2,
\]
其中 $\{z;w\}$ 是 Schwarz 导数，定义为
\[
  \{z;w\}=\frac{d^3z}{dw^3}\Big/\frac{dz}{dw}
  -\frac{3}{2}\left(\frac{d^2z}{dw^2}\Big/\frac{dz}{dw}\right)^2 .
\]
\end{theorem}

